\documentclass[11pt]{article}
% ===========================================================================
%  Shared preamble for the AP Calculus AB exam-preparation worksheets.
%  Input right after \documentclass. Files beginning with "_" are not compiled
%  on their own. Navy and gold, to match the course's lesson pages.
% ===========================================================================
\usepackage[letterpaper,top=2.3cm,bottom=2.2cm,left=1.7cm,right=1.7cm,headheight=28pt,footskip=30pt]{geometry}
\usepackage{amsmath,amssymb}
\usepackage[T1]{fontenc}
\usepackage[utf8]{inputenc}
\usepackage{lmodern}
\usepackage{xcolor}
\usepackage[most]{tcolorbox}
\usepackage{enumitem}
\usepackage{fancyhdr}
\usepackage{array}
\usepackage{tikz}
\usepackage{pgfplots}
\pgfplotsset{compat=1.18}

\definecolor{navy}{HTML}{161E54}
\definecolor{ink}{HTML}{27357E}
\definecolor{gold}{HTML}{B98A00}
\definecolor{goldbg}{HTML}{FFFDF3}
\definecolor{inkbg}{HTML}{F6F7FD}
\definecolor{rule}{HTML}{C9CFF0}
\definecolor{okgreen}{HTML}{1B7A44}
\definecolor{okbg}{HTML}{F1F8F3}

\pagestyle{fancy}
\fancyhf{}
\fancyhead[L]{\footnotesize Name: \underline{\hspace{4cm}}}
\fancyhead[C]{\textbf{\color{navy}Integration Academy}\\[-2pt]{\scriptsize Calculus AB \textperiodcentered{} AP\textsuperscript{\textregistered} Exam Preparation}}
\fancyfoot[L]{\tiny\color{black!55}\parbox[t]{0.78\textwidth}{AP\textsuperscript{\textregistered} is a registered trademark of the College Board, which was not involved in the production of, and does not endorse, this material. \textcopyright{} Integration Academy.}}
\fancyfoot[R]{\scriptsize\thepage}
\renewcommand{\headrulewidth}{1.2pt}
\renewcommand{\headrule}{\hbox to\headwidth{\color{gold}\leaders\hrule height \headrulewidth\hfill}}
\renewcommand{\footrulewidth}{0pt}
\setlength{\parindent}{0pt}
\setlength{\parskip}{4pt}

% Title block: \wstitle{code}{title}{unit line}
\newcommand{\wstitle}[3]{\begin{center}{\small\bfseries\color{gold}WORKSHEET #1}\\[3pt]{\LARGE\bfseries\color{navy}#2}\\[3pt]{\small\color{black!60}#3}\end{center}\vspace{2pt}}
% Section heading
\newcommand{\sech}[1]{\par\vspace{10pt}{\large\bfseries\color{navy}#1}\par\nobreak\vspace{2pt}{\color{rule}\hrule height0.8pt}\vspace{7pt}}
% A level tag printed beside a group of questions
\newcommand{\level}[1]{\par\vspace{6pt}{\footnotesize\bfseries\color{gold}\MakeUppercase{#1}}\par\nobreak\vspace{2pt}}

\newtcolorbox{ideabox}{colframe=ink,colback=inkbg,coltitle=white,colbacktitle=ink,fonttitle=\bfseries\small,title={Key ideas},arc=2pt,boxrule=0.8pt,left=8pt,right=8pt}
\newtcolorbox{example}[1]{colframe=rule,colback=white,coltitle=navy,colbacktitle=inkbg,fonttitle=\bfseries\small,title={#1},arc=2pt,boxrule=0.8pt,borderline west={3pt}{0pt}{ink},left=10pt,right=8pt}
\newtcolorbox{frbox}[1]{colframe=gold,colback=goldbg,coltitle=white,colbacktitle=gold,fonttitle=\bfseries\small,title={#1},arc=2pt,boxrule=0.8pt,breakable,left=8pt,right=8pt}
\newtcolorbox{solbox}{colframe=okgreen,colback=okbg,coltitle=white,colbacktitle=okgreen,fonttitle=\bfseries\small,title={Worked solutions},arc=2pt,boxrule=0.8pt,breakable,left=8pt,right=8pt}
\newtcolorbox{scorebox}[1]{colframe=navy,colback=white,coltitle=white,colbacktitle=navy,fonttitle=\bfseries\small,title={#1},arc=2pt,boxrule=0.8pt,breakable,left=8pt,right=8pt}

\newcommand{\soln}{\par\smallskip\textbf{\color{ink}Solution.}\ }
\newcommand{\result}[1]{\par\smallskip\textbf{#1}}
% blank working space of a given height
\newcommand{\work}[1]{\par\vspace{#1}}
% a point value in the margin of a scoring guide
\newcommand{\pts}[1]{\hfill{\footnotesize\bfseries\color{gold}[#1]}}

% Coordinate grid for graph questions: \begin{apgraph}{xmin}{xmax}{ymin}{ymax} ... \end{apgraph}
\newenvironment{apgraph}[4]{\begin{center}\begin{tikzpicture}\begin{axis}[
  axis lines=middle, xmin=#1, xmax=#2, ymin=#3, ymax=#4,
  grid=both, grid style={black!12}, xtick distance=1, ytick distance=1,
  ticklabel style={font=\scriptsize}, xlabel={$x$}, ylabel={$y$},
  width=10cm, height=7.2cm, samples=80, clip=false,
  every axis plot/.append style={line width=1.1pt, color=ink}]}{\end{axis}\end{tikzpicture}\end{center}}
\newcommand{\closeddot}[2]{\addplot[only marks, mark=*, mark size=2.2pt, color=ink] coordinates {(#1,#2)};}
\newcommand{\opendot}[2]{\addplot[only marks, mark=*, mark size=2.2pt, color=ink, mark options={fill=white}] coordinates {(#1,#2)};}

% Multiple-choice options: four across, or stacked when the options are long.
\newcommand{\opts}[4]{\par\vspace{3pt}\begin{tabular}{@{}p{0.235\linewidth}p{0.235\linewidth}p{0.235\linewidth}p{0.235\linewidth}@{}}(A)\ #1 & (B)\ #2 & (C)\ #3 & (D)\ #4\end{tabular}\par\vspace{9pt}}
\newcommand{\optsv}[4]{\par\vspace{3pt}\begin{tabular}{@{}l@{\ \ }p{0.9\linewidth}@{}}(A) & #1\\(B) & #2\\(C) & #3\\(D) & #4\end{tabular}\par\vspace{9pt}}

\fancyhead[R]{\footnotesize\color{black!60}Worksheet 9.8}
\begin{document}

\wstitle{9.8}{Multiple-Choice Technique and the Last Two Weeks}{Unit 9: Exam Strategy and Free-Response Types}

\begin{ideabox}
\begin{itemize}[leftmargin=*,itemsep=2pt,topsep=2pt]
  \item Eliminate first. A sign, a size or a unit often removes two options before any calculus.
  \item Do only the work the question needs. If the options differ in one feature, find that feature.
  \item Check an antiderivative by differentiating it, and a formula by trying an easy value.
  \item Two minutes without progress: choose, mark the question, move on. Never leave a blank.
\end{itemize}
\end{ideabox}

\sech{Worked examples}
\begin{example}{Example 1 \textperiodcentered{} Reading the leading terms}
$\displaystyle\lim_{x\to\infty}\dfrac{3x^2+1}{5x^2-x}$ is \\[2pt] (A) $0$ \qquad (B) $\dfrac35$ \qquad (C) $1$ \qquad (D) $\infty$\soln Equal degrees: the ratio of the leading coefficients.\result{(B).}
\end{example}

\begin{example}{Example 2 \textperiodcentered{} Using symmetry}
$\displaystyle\int_{-2}^{2}(x^3+1)\,dx$ is \\[2pt] (A) $0$ \qquad (B) $2$ \qquad (C) $4$ \qquad (D) $8$\soln The odd part $x^3$ integrates to 0 on a symmetric interval, leaving $\displaystyle\int_{-2}^{2}1\,dx$.\result{(C).}
\end{example}

\begin{example}{Example 3 \textperiodcentered{} Checking by differentiating}
An antiderivative of $x\cos\!\left(x^2\right)$ is \\[2pt] (A) $\sin\!\left(x^2\right)$ \qquad (B) $\tfrac12\sin\!\left(x^2\right)$ \qquad (C) $2\sin\!\left(x^2\right)$ \qquad (D) $-\tfrac12\sin\!\left(x^2\right)$\soln Differentiate each option: $\tfrac12\sin\!\left(x^2\right)$ gives $x\cos\!\left(x^2\right)$.\result{(B).}
\end{example}

\newpage
\sech{Practice}
\emph{No calculator on this page.}
\level{Core}
\begin{enumerate}[leftmargin=*,itemsep=2pt]
  \item If $f(x)=\sin(3x)$, then $f'(0)$ is \\[2pt] (A) $0$ \qquad (B) $1$ \qquad (C) $3$ \qquad (D) $-3$\work{2.3cm}
  \item $\displaystyle\int_{0}^{1}4x^3\,dx$ is \\[2pt] (A) $1$ \qquad (B) $4$ \qquad (C) $12$ \qquad (D) $\tfrac14$\work{2.3cm}
  \item If $f'(x)<0$ and $f''(x)>0$, the graph of $f$ is \\[2pt] (A) increasing, concave up \qquad (B) decreasing, concave up \qquad (C) decreasing, concave down \qquad (D) increasing, concave down\work{2.3cm}
  \item $\displaystyle\lim_{x\to0}\dfrac{\sin(5x)}{x}$ is \\[2pt] (A) $0$ \qquad (B) $1$ \qquad (C) $5$ \qquad (D) nonexistent\work{2.3cm}
\end{enumerate}

\level{Exam level}
\begin{enumerate}[leftmargin=*,itemsep=2pt,start=5]
  \item $\dfrac{d}{dx}\displaystyle\int_{1}^{x^2}\dfrac{1}{1+t^4}\,dt$ is \\[2pt] (A) $\dfrac{1}{1+x^8}$ \qquad (B) $\dfrac{2x}{1+x^8}$ \qquad (C) $\dfrac{2x}{1+x^4}$ \qquad (D) $\dfrac{1}{1+x^4}$\work{2.8cm}
  \item If $f'(x)=(x-1)^2(x-3)$, then $f$ has \\[2pt] (A) a relative minimum at $x=1$ \qquad (B) a relative maximum at $x=1$ \qquad (C) a relative minimum at $x=3$ \qquad (D) a relative maximum at $x=3$\work{2.8cm}
  \item The average value of $3x^2$ on $[0,2]$ is \\[2pt] (A) $4$ \qquad (B) $8$ \qquad (C) $12$ \qquad (D) $2$\work{2.8cm}
  \item The line tangent to $y=e^{2x}$ at $x=0$ is used to approximate $e^{0.2}$. The approximation is \\[2pt] (A) $1.02$ \qquad (B) $1.1$ \qquad (C) $1.2$ \qquad (D) $1.4$\work{2.8cm}
\end{enumerate}

\begin{frbox}{Question 9 \textperiodcentered{} Free-response practice}
Answer each part in one or two sentences.
\begin{enumerate}[label=(\alph*),leftmargin=*,itemsep=2pt]
  \item Without evaluating it, explain why $\displaystyle\int_{0}^{2}(x^2+1)\,dx$ cannot equal 1.\work{2.6cm}
  \item A question offers $xe^{x}+e^{x}$ and $xe^{x}$ as possible derivatives of $xe^{x}$. Decide between them by considering the slope of the graph at $x=0$, without the product rule.\work{2.6cm}
  \item You have used 30 of the 38 minutes of the calculator multiple-choice part and answered 9 of 13 questions. Describe what you should do with the remaining time.\work{2.6cm}
\end{enumerate}
\end{frbox}

\newpage
\begin{solbox}
\begin{enumerate}[leftmargin=*,itemsep=4pt]
  \item (C): $f'(x)=3\cos(3x)$.
  \item (A): $\Big[x^4\Big]_0^1=1$.
  \item (B).
  \item (C).
  \item (B): substitute the upper limit and multiply by its derivative, $2x$.
  \item (C): $f'$ changes sign, from negative to positive, only at $x=3$.
  \item (A): $\dfrac12\Big[x^3\Big]_0^2=4$.
  \item (C): the tangent line is $y=1+2x$, and $e^{0.2}=e^{2(0.1)}$, so use $x=0.1$.
  \item \begin{enumerate}[label=(\alph*),leftmargin=*,itemsep=3pt]
    \item On $[0,2]$ the integrand is at least 1, so the integral is at least $1\times2=2$.
    \item Near $x=0$, $e^{x}\approx1$, so $xe^{x}\approx x$ and the slope at $x=0$ is 1. The first option gives $0+1=1$ and the second gives 0, so the first is correct.
    \item About 2 minutes for each of the 4 remaining questions. Answer first those that need one calculator step, choose an option for every question, and leave none blank.
  \end{enumerate}
\end{enumerate}
\end{solbox}

\end{document}
